\documentclass[11pt,a4paper]{article}

% ============================================================================
% PACKAGES
% ============================================================================
\usepackage[utf8]{inputenc}
\usepackage[T1]{fontenc}
\usepackage{amsmath}
\usepackage{amsfonts}
\usepackage{amstext}
\usepackage{graphicx}
\usepackage{hyperref}
\usepackage{listings}
\usepackage{xcolor}
\usepackage{booktabs}
\usepackage{geometry}
\usepackage{fancyhdr}
\usepackage{titlesec}
\usepackage{enumitem}
\usepackage{algorithm}
\usepackage{algpseudocode}
\usepackage{amsopn}
\usepackage{float}
\usepackage{xurl}
\usepackage{caption}
\usepackage{subcaption}

% ============================================================================
% DOCUMENT SETTINGS
% ============================================================================
\geometry{margin=1in}
\hypersetup{
    colorlinks=true,
    linkcolor=blue,
    filecolor=magenta,
    urlcolor=cyan,
    citecolor=blue
}

% Code listing settings
\lstset{
    basicstyle=\ttfamily\small,
    breaklines=true,
    frame=single,
    language=Python,
    keywordstyle=\color{blue}\bfseries,
    stringstyle=\color{red},
    commentstyle=\color{gray}\itshape,
    numbers=left,
    numberstyle=\tiny\color{gray},
    backgroundcolor=\color{gray!5},
    showstringspaces=false,
    tabsize=4,
    captionpos=b
}

% ============================================================================
% TITLE
% ============================================================================
\title{
    \textbf{Search and Rescue Multi-Agent Simulation}\\[0.5em]
    \Large Technical Documentation\\[1em]
    \normalsize ELTE Collective Intelligence Course --- Assignment 2
}
\author{
    Adorján Nagy-Mohos\\
    \texttt{d5vd5e@inf.elte.hu}
    \and
    Máté Kovács\\
    \texttt{u5bky4@inf.elte.hu}
    \and
    Sándor Baranyi\\
    \texttt{ct9xfj@inf.elte.hu}
}
\date{\today}

% ============================================================================
% DOCUMENT
% ============================================================================
\begin{document}

\maketitle

\begin{abstract}
This document provides comprehensive technical documentation for the Search and Rescue Multi-Agent Simulation project.
The simulation utilizes the PettingZoo library for modeling a cooperative multi-agent environment where intelligent rescuer agents must navigate obstacles, locate victims, and guide them to designated safe zones.
The project implements Multi-Agent Proximal Policy Optimization (MAPPO) using the TorchRL framework for training decentralized policies with centralized value estimation.
This documentation covers the environment design, agent dynamics, reward mechanisms, training pipeline, and evaluation methodology.
\end{abstract}

\tableofcontents
\newpage

% ============================================================================
\section{Introduction}
\label{sec:introduction}
% ============================================================================

This document provides an in-depth overview of the \textbf{Search and Rescue Multi-Agent Simulation}, which utilizes the PettingZoo library for modeling a cooperative multi-agent environment.
The simulation is designed to replicate dynamic rescue operations where rescuers navigate an environment containing obstacles, victims, and predefined safe zones to achieve specific objectives.

The primary goal of this project is to develop and evaluate multi-agent reinforcement learning algorithms in a challenging cooperative setting that requires:
\begin{itemize}
    \item \textbf{Coordination}: Multiple agents must work together efficiently without explicit communication.
    \item \textbf{Partial Observability}: Agents have limited vision and must navigate with incomplete information.
    \item \textbf{Credit Assignment}: Individual agents must learn their contribution to team success.
    \item \textbf{Dynamic Interactions}: Victims respond to agent behavior through a commitment mechanism.
\end{itemize}

The following chapters describe the functionalities, components, mathematical formulations, and implementation details of the project.

% ============================================================================
\section{Functionalities}
\label{sec:functionalities}
% ============================================================================

This project presents a comprehensive simulation of a Search and Rescue scenario, designed within a multi-agent environment using the PettingZoo library.
The environment models a dynamic rescue operation where cooperative agents must navigate a space containing both obstacles and targets to achieve predefined goals.
The simulation incorporates the following components:

% ----------------------------------------------------------------------------
\subsection{Environment Setup}
\label{subsec:env_setup}
% ----------------------------------------------------------------------------

The environment is configured through a set of parameters that define the simulation characteristics:

\begin{itemize}
    \item \textbf{Configurable Parameters}: Users can define key elements of the environment, including:
    \begin{itemize}
        \item Number of rescuers (agents responsible for guiding victims).
        \item Number of victims (missing agents to be rescued).
        \item Number of trees (obstacles obstructing navigation).
        \item Number of safe zones (designated target areas for rescue).
    \end{itemize}

    \item \textbf{Agent Interaction}: Rescuers interact with victims using a commitment-based following mechanism.
    Upon approaching a victim within the follow radius, the victim commits to following that rescuer.
    This commitment system implements hysteresis to prevent rapid switching between rescuers, providing stable victim-agent relationships.

    \item \textbf{Safe Zone Placement}: Safe zones are fixed at the four corners of the environment.
    Each zone is assigned a specific type (e.g., A, B, C, D), and only victims of the matching type are considered successfully rescued when they reach the zone.
    This type-matching requirement introduces an additional layer of strategic planning for the agents.
\end{itemize}

% ----------------------------------------------------------------------------
\subsection{Rescuers (Agents)}
\label{subsec:rescuers}
% ----------------------------------------------------------------------------

Rescuers are intelligent agents that serve as the primary actors in the simulation.
They exhibit the following characteristics:

\begin{itemize}
    \item \textbf{Motion Constraints}: Operate under physical constraints such as limited speed, acceleration, and maneuverability.
    The maximum speed is capped at 0.08 units per timestep, and actions are applied through a velocity-based control scheme.

    \item \textbf{Reactive Behavior}: Demonstrate dynamic behavior by interacting with victims and avoiding obstacles.
    Agents must adapt their trajectories in real-time based on changing environmental conditions.

    \item \textbf{Commitment Mechanism}: Use a commitment-based system to guide victims.
    When a rescuer approaches a victim within the follow radius, the victim becomes committed to following that rescuer until either rescued or the rescuer moves too far away.

    \item \textbf{Collision Response}: Agents implement collision detection and response mechanisms for walls, trees, and other agents.
    Tree collisions result in velocity reflection with damping, while agent-agent proximity triggers soft repulsion forces.
\end{itemize}

% ----------------------------------------------------------------------------
\subsection{Victims (Missing Agents)}
\label{subsec:victims}
% ----------------------------------------------------------------------------

Victims represent passive entities scattered within the environment that require rescue.
Their characteristics include:

\begin{itemize}
    \item \textbf{Type Assignment}: Each victim is assigned a specific type (A, B, C, or D) that corresponds to safe zone types.
    Types are assigned cyclically: victim $i$ receives type $i \mod 4$.

    \item \textbf{Matching Requirement}: Victims must be guided to their matching safe zones to be considered rescued.
    A victim of type A can only be saved at the type A safe zone.

    \item \textbf{Following Behavior}: When committed to a rescuer, victims follow the assigned agent with a following force of 0.04 units.
    The victim velocity is updated as shown in Equation~\ref{eq:victim_velocity}.
    \begin{equation}
    \label{eq:victim_velocity}
    \mathbf{v}_{\text{victim}}^{t+1} = 0.8 \cdot \mathbf{v}_{\text{victim}}^{t} + 0.04 \cdot \hat{\mathbf{d}}_{\text{to\_rescuer}}
    \end{equation}

    \item \textbf{Brownian Motion}: When not committed to any rescuer, victims exhibit simple Brownian motion with noise magnitude 0.0075, simulating random wandering behavior.

    \item \textbf{Commitment Hysteresis}: To prevent rapid switching between rescuers, the commitment system implements hysteresis.
    A victim releases its current assignment only when the assigned rescuer moves beyond 1.5 times the follow radius, and switches to a new rescuer only if the new rescuer is significantly closer (less than 0.6 times the current distance).
\end{itemize}

% ----------------------------------------------------------------------------
\subsection{Safe Zones}
\label{subsec:safe_zones}
% ----------------------------------------------------------------------------

Safe zones serve as stationary rescue targets for victims and define the success criteria:

\begin{itemize}
    \item \textbf{Corner Placement}: Safe zones are predefined at the four corners of the environment at positions:
    \begin{itemize}
        \item Type 0 (A): Top-left $(-0.9, 0.9)$
        \item Type 1 (B): Top-right $(0.9, 0.9)$
        \item Type 2 (C): Bottom-left $(-0.9, -0.9)$
        \item Type 3 (D): Bottom-right $(0.9, -0.9)$
    \end{itemize}

    \item \textbf{Type Identification}: Each safe zone is assigned a unique type, represented by distinct colors (red, green, blue, yellow) for visual identification during rendering.

    \item \textbf{Rescue Condition}: Successful rescues occur only when victims of matching types reach their corresponding safe zones.
    The rescue is triggered when the Euclidean distance between victim and safe zone falls below the rescue radius of 0.15 units.

    \item \textbf{Optional Randomization}: The environment supports optional randomization of safe zone types while maintaining corner positions, controlled by the \texttt{randomize\_safe\_zones} parameter.
\end{itemize}

% ----------------------------------------------------------------------------
\subsection{Obstacles (Trees)}
\label{subsec:obstacles}
% ----------------------------------------------------------------------------

Trees are static obstacles that add complexity to agent navigation:

\begin{itemize}
    \item \textbf{Random Distribution}: Trees are randomly distributed within the environment during initialization and maintain fixed positions throughout an episode.
    The random placement uses uniform sampling within the range $[-0.5, 0.5]$ to avoid spawning too close to corners.

    \item \textbf{Collision Penalty}: Agents are penalized for colliding with trees.
    Collisions are detected using Euclidean distance between agent and tree centers, with collision occurring when the condition in Equation~\ref{eq:collision} is satisfied.
    \begin{equation}
    \label{eq:collision}
    \|\mathbf{x}_{\text{agent}} - \mathbf{x}_{\text{tree}}\| < r_{\text{agent}} + r_{\text{tree}}
    \end{equation}
    where $r_{\text{agent}} = 0.03$ and $r_{\text{tree}} = 0.05$.

    \item \textbf{Path Obstruction}: Obstacles add complexity by blocking direct paths between rescuers and victims or safe zones.
    This requires agents to develop navigation strategies that avoid obstacles while pursuing objectives.

    \item \textbf{Line-of-Sight Blocking}: Trees obstruct the line-of-sight between agents and other entities.
    When an entity is occluded by a tree, it is masked in the observation with zeros, introducing partial observability into the environment.
\end{itemize}

% ----------------------------------------------------------------------------
\subsection{Collision Avoidance}
\label{subsec:collision_avoidance}
% ----------------------------------------------------------------------------

The environment implements comprehensive collision avoidance mechanisms:

\begin{itemize}
    \item \textbf{Agent-Agent Avoidance}: Rescuers apply soft repulsion forces when within proximity (0.15 units) of each other.
    The repulsion force magnitude is given by Equation~\ref{eq:repulsion}.
    \begin{equation}
    \label{eq:repulsion}
    F_{\text{repulsion}} = 0.005 \cdot \frac{r_{\text{threshold}} - d}{d}
    \end{equation}
    where $d$ is the distance between agents.
    This prevents clustering while maintaining smooth agent motion.

    \item \textbf{Obstacle Avoidance}: Agents are penalized ($-1.0$ reward) for colliding with trees.
    Upon collision, the agent's velocity is reflected about the collision normal with damping factor of approximately 0.5, as shown in Equation~\ref{eq:reflection}.
    \begin{equation}
    \label{eq:reflection}
    \mathbf{v}' = \mathbf{v} - 1.5(\mathbf{v} \cdot \hat{\mathbf{n}})\hat{\mathbf{n}}
    \end{equation}

    \item \textbf{Wall Collision Handling}: When agents cross world boundaries $[-1, 1]$, their position is clamped to the boundary and the corresponding velocity component is inverted with 0.5 damping, as shown in Equation~\ref{eq:wall}.
    \begin{equation}
    \label{eq:wall}
    v_{\text{axis}} \leftarrow -0.5 \cdot v_{\text{axis}}
    \end{equation}

    \item \textbf{Line-of-Sight Blocking}: Trees obstruct vision between agents.
    The occlusion check determines whether a tree's bounding region intersects the line segment between observer and target.
    When visibility is blocked, agents must navigate indirectly.

    \item \textbf{Reactive Behavior}: The environment promotes indirect navigation strategies, such as moving tangentially to obstacles, to simulate realistic path planning behaviors.
\end{itemize}

% ----------------------------------------------------------------------------
\subsection{Multi-Agent Interaction and Dynamics}
\label{subsec:multi_agent}
% ----------------------------------------------------------------------------

The multi-agent aspects of the environment create complex emergent behaviors:

\begin{itemize}
    \item \textbf{Agent Collaboration}: Although agents operate independently with decentralized policies, their shared objective of rescuing victims promotes emergent collaborative behaviors, such as efficient task distribution and implicit role assignment.

    \item \textbf{Motion Dynamics}: Agents operate under physical constraints including speed limits and acceleration.
    The velocity update equation is given by Equation~\ref{eq:velocity_update}.
    \begin{equation}
    \label{eq:velocity_update}
    \mathbf{v}^{t+1} = 0.8 \cdot \mathbf{v}^t + 0.1 \cdot \mathbf{a}^t
    \end{equation}
    where $\mathbf{a}^t \in [-1, 1]^2$ is the action (acceleration) and velocity is clipped to maximum speed 0.08.

    \item \textbf{Credit Assignment Challenge}: With multiple agents contributing to team success, the environment presents a credit assignment problem.
    The reward structure is designed to attribute rewards to individual agents based on their specific contributions (e.g., the assigned rescuer receives the rescue bonus).

    \item \textbf{Homogeneous Agents}: All rescuer agents share the same capabilities and policy parameters, promoting the learning of general strategies that work across different agent positions and roles.
\end{itemize}

% ----------------------------------------------------------------------------
\subsection{Training and Evaluation Framework}
\label{subsec:training_eval}
% ----------------------------------------------------------------------------

The project provides comprehensive training and evaluation capabilities:

\begin{itemize}
    \item \textbf{Training Framework}: The simulation uses TorchRL's implementation of Multi-Agent PPO (MAPPO) for training.
    Key components include:
    \begin{itemize}
        \item SyncDataCollector for parallel experience collection.
        \item ClipPPOLoss with configurable clipping threshold.
        \item Generalized Advantage Estimation (GAE) for advantage computation.
        \item Adam optimizer with configurable learning rate.
    \end{itemize}

    \item \textbf{Curriculum Learning}: Optional curriculum learning support allows for progressive difficulty adjustment:
    \begin{itemize}
        \item Training starts with fewer obstacles (configurable minimum).
        \item Difficulty gradually increases through defined stages.
        \item Each stage maintains consistent observation space size.
    \end{itemize}

    \item \textbf{Configurable Timesteps}: Users can adjust the training duration through the \texttt{total\_timesteps} parameter to control convergence and computational requirements.

    \item \textbf{Performance Logging}: Comprehensive logging is provided through:
    \begin{itemize}
        \item \textbf{TensorBoard}: Provides real-time performance visualization including loss curves, reward trajectories, and training statistics.
        \item \textbf{Checkpoint Saving}: Model weights and configuration are saved for later evaluation and fine-tuning.
    \end{itemize}

    \item \textbf{Evaluation Pipeline}: The evaluation system supports:
    \begin{itemize}
        \item Manual model selection or automatic detection of latest checkpoint.
        \item Environment configuration loading from saved checkpoints.
        \item Per-episode metrics collection and aggregation.
        \item Real-time rendering for visual verification.
    \end{itemize}
\end{itemize}

% ----------------------------------------------------------------------------
\subsection{Reward System}
\label{subsec:reward_system}
% ----------------------------------------------------------------------------

The reward system governs agent behavior through a combination of positive rewards and penalties.

\subsubsection{Positive Rewards}
\label{subsubsec:positive_rewards}

\begin{itemize}
    \item \textbf{Successful Rescue}: A large reward of $+100.0$ is awarded to the rescuer who was assigned to (escorting) the victim when it reaches its matching safe zone.
    This is formalized in Equation~\ref{eq:success_reward}.
    \begin{equation}
    \label{eq:success_reward}
    R_{\text{success}} = \begin{cases}
        100.0 & \text{if victim reaches matching safe zone} \\
        0 & \text{otherwise}
    \end{cases}
    \end{equation}

    \item \textbf{Escort Reward}: The assigned rescuer receives a continuous reward proportional to how close the victim is to its target safe zone, as shown in Equation~\ref{eq:escort_reward}.
    \begin{equation}
    \label{eq:escort_reward}
    R_{\text{escort}} = \min\left(\frac{1.0}{d_{\text{victim-safezone}} + 0.1}, 10.0\right)
    \end{equation}
    This reward is capped at 10.0 to prevent infinite rewards when very close.

    \item \textbf{Distance Shaping}: All rescuers receive a shaping reward for reducing distance to unsaved victims, as shown in Equation~\ref{eq:shaping_reward}.
    \begin{equation}
    \label{eq:shaping_reward}
    R_{\text{shaping}} = 0.1 \cdot (d_{\text{prev}} - d_{\text{current}})
    \end{equation}
    where positive values indicate the agent moved closer to victims.
\end{itemize}

\subsubsection{Penalties}
\label{subsubsec:penalties}

\begin{itemize}
    \item \textbf{Low Velocity Penalty}: Rescuers with speed below 0.01 receive a penalty of $-0.05$ per step to encourage active exploration.
    This is formalized in Equation~\ref{eq:velocity_penalty}.
    \begin{equation}
    \label{eq:velocity_penalty}
    R_{\text{velocity}} = \begin{cases}
        -0.05 & \text{if } \|\mathbf{v}\| < 0.01 \\
        0 & \text{otherwise}
    \end{cases}
    \end{equation}

    \item \textbf{Tree Collision}: A penalty of $-1.0$ is applied when an agent collides with a tree obstacle.

    \item \textbf{Boundary Violation}: A penalty of $-1.0$ is applied when an agent's position exceeds 0.95 in any dimension.

    \item \textbf{Agent Collision}: A stronger penalty of $-5.0$ is applied to both agents when they are within collision distance (0.15 units) to discourage clustering.
\end{itemize}

\subsubsection{Reward Summary Table}
\label{subsubsec:reward_table}

\begin{table}[H]
\centering
\begin{tabular}{lll}
\toprule
\textbf{Event} & \textbf{Reward} & \textbf{Recipient} \\
\midrule
Successful rescue & $+100.0$ & Assigned rescuer \\
Escorting victim & $+1.0 / (d + 0.1)$ & Assigned rescuer (per step) \\
Distance shaping & $+0.1 \times \Delta d$ & All rescuers (getting closer) \\
Low velocity & $-0.05$ & Stationary rescuer (per step) \\
Tree collision & $-1.0$ & Colliding rescuer \\
Boundary violation & $-1.0$ & Violating rescuer \\
Agent collision & $-5.0$ & Both colliding rescuers \\
\bottomrule
\end{tabular}
\caption{Complete reward structure for the Search and Rescue environment.}
\label{tab:rewards}
\end{table}

\subsubsection{Combined Reward Function}
\label{subsubsec:combined_reward}

The total reward for a rescuer agent $i$ at timestep $t$ is given by Equation~\ref{eq:total_reward}.
\begin{equation}
\label{eq:total_reward}
R_i^t = R_{\text{success}} + R_{\text{escort}} + R_{\text{shaping}} + R_{\text{velocity}} - R_{\text{penalties}}
\end{equation}
where $R_{\text{penalties}}$ aggregates collision and boundary penalties.

% ----------------------------------------------------------------------------
\subsection{Observation Space}
\label{subsec:observation_space}
% ----------------------------------------------------------------------------

Each rescuer agent receives a structured observation vector that provides local information about the environment.
The observation components are described below.

\subsubsection{Observation Components}
\label{subsubsec:obs_components}

\begin{enumerate}
    \item \textbf{Self State} (4 values):
    \begin{itemize}
        \item Velocity: $[v_x, v_y]$
        \item Position: $[x, y]$
    \end{itemize}

    \item \textbf{Agent ID} ($N_{\text{rescuers}}$ values):
    \begin{itemize}
        \item One-hot encoding of agent index for symmetry breaking.
        \item Allows shared policy to differentiate between agents.
    \end{itemize}

    \item \textbf{Safe Zones} ($N_{\text{safe\_zones}} \times 3$ values):
    \begin{itemize}
        \item Relative position: $[x_{\text{rel}}, y_{\text{rel}}]$
        \item Type index: normalized type value.
    \end{itemize}

    \item \textbf{Trees} ($N_{\text{trees}} \times 2$ values):
    \begin{itemize}
        \item Relative position: $[x_{\text{rel}}, y_{\text{rel}}]$
        \item Masked as $[0, 0]$ if tree is not within vision radius.
    \end{itemize}

    \item \textbf{Victims} ($N_{\text{victims}} \times 3$ values):
    \begin{itemize}
        \item Relative position: $[x_{\text{rel}}, y_{\text{rel}}]$
        \item Type index (masked as $[0, 0, -1]$ if not visible or already saved).
    \end{itemize}
\end{enumerate}

\subsubsection{Observation Dimension}
\label{subsubsec:obs_dimension}

The total observation dimension is given by Equation~\ref{eq:obs_dim}.
\begin{equation}
\label{eq:obs_dim}
D_{\text{obs}} = 4 + N_{\text{rescuers}} + 3 \times N_{\text{safe\_zones}} + 2 \times N_{\text{trees}} + 3 \times N_{\text{victims}}
\end{equation}

\subsubsection{Visibility and Masking}
\label{subsubsec:visibility}

Observations are partially masked based on visibility constraints:
\begin{itemize}
    \item Entities beyond the vision radius are masked with zeros.
    \item Entities occluded by trees (line-of-sight blocking) are masked.
    \item Saved victims are masked to indicate they no longer need rescue.
\end{itemize}

The mathematical observation for agent $i$ is given by Equation~\ref{eq:observation}.
\begin{equation}
\label{eq:observation}
\mathbf{o}_i = [\mathbf{v}_i, \mathbf{x}_i, \mathbf{e}_i^{\text{id}}, \mathbf{x}_{\text{rel\_safezones}}, \mathbf{x}_{\text{rel\_trees}}, \mathbf{x}_{\text{rel\_victims}}]
\end{equation}
where $\mathbf{x}_{\text{rel}}$ denotes relative positions computed as shown in Equation~\ref{eq:relative_pos}.
\begin{equation}
\label{eq:relative_pos}
\mathbf{x}_{\text{rel}} = \mathbf{x}_{\text{entity}} - \mathbf{x}_i
\end{equation}

% ============================================================================
\section{Structure of the Code}
\label{sec:code_structure}
% ============================================================================

The project is organized into modular components, each responsible for a specific aspect of the simulation pipeline.
This section provides detailed documentation of each module.

% ----------------------------------------------------------------------------
\subsection{Custom Environment: sar\_env.py}
\label{subsec:sar_env}
% ----------------------------------------------------------------------------

The \texttt{sar\_env.py} file defines the multi-agent search-and-rescue environment using the PettingZoo ParallelEnv API.
It models agents (rescuers and victims), landmarks (trees and safe zones), and their interactions.

\subsubsection{Environment Dynamics}
\label{subsubsec:env_dynamics}

The environment operates in a 2D continuous space where each agent $i$ has a position $\mathbf{x}_i$ and velocity $\mathbf{v}_i$.
These are defined in Equation~\ref{eq:agent_state}.
\begin{equation}
\label{eq:agent_state}
\mathbf{x}_i = [x_i, y_i], \quad \mathbf{v}_i = [v_{x_i}, v_{y_i}]
\end{equation}

Positions are updated according to velocity as shown in Equation~\ref{eq:position_update}.
\begin{equation}
\label{eq:position_update}
\mathbf{x}_i(t+1) = \mathbf{x}_i(t) + \mathbf{v}_i(t)
\end{equation}

The world boundaries are defined as shown in Equation~\ref{eq:boundaries}.
\begin{equation}
\label{eq:boundaries}
x_{\min} \leq x_i \leq x_{\max}, \quad y_{\min} \leq y_i \leq y_{\max}
\end{equation}
where $x_{\min} = y_{\min} = -1$ and $x_{\max} = y_{\max} = 1$.

Agents are penalized and their velocity is reflected if they move out of bounds.

\subsubsection{Collision Detection}
\label{subsubsec:collision_detection}

Collisions occur when the Euclidean distance between two entities is less than the sum of their radii, as shown in Equation~\ref{eq:collision_general}.
\begin{equation}
\label{eq:collision_general}
\|\mathbf{x}_1 - \mathbf{x}_2\| \leq r_1 + r_2
\end{equation}
where $r_1$ and $r_2$ are the radii of the entities.

\begin{lstlisting}[caption={Collision detection implementation.},label={lst:collision}]
def is_collision(self, pos1, pos2, r1, r2):
    delta_pos = pos1 - pos2
    dist = np.sqrt(np.sum(np.square(delta_pos)))
    dist_min = r1 + r2
    return dist < dist_min
\end{lstlisting}

\subsubsection{Reward Mechanism Implementation}
\label{subsubsec:reward_impl}

The reward computation is implemented in the \texttt{\_compute\_rewards} method.

\begin{lstlisting}[caption={Reward computation pseudocode.},label={lst:rewards}]
def _compute_rewards(self, rewards):
    # Distance shaping
    for agent in agents:
        delta = prev_dist[agent] - current_dist[agent]
        rewards[agent] += delta * 0.1

    # Low velocity penalty
    for agent in agents:
        if velocity[agent] < 0.01:
            rewards[agent] -= 0.05

    # Check rescues
    for victim in victims:
        if not saved and dist_to_matching_zone < rescue_radius:
            saved[victim] = True
            rewards[assigned_rescuer] += 100.0

    # Escort rewards for assigned rescuers
    for victim in unsaved_victims:
        if has_assignment:
            proximity_reward = min(1.0 / (dist_to_zone + 0.1), 10.0)
            rewards[assigned_agent] += proximity_reward

    # Boundary penalties
    for agent in agents:
        if abs(pos) > 0.95:
            rewards[agent] -= 1.0

    # Agent collision penalties
    for pair in agent_pairs:
        if dist < 0.15:
            rewards[agent_i] -= 5.0
            rewards[agent_j] -= 5.0
\end{lstlisting}

\subsubsection{Observation Space Construction}
\label{subsubsec:obs_construction}

The observation for agent $i$ is constructed by concatenating the relevant components.

\begin{lstlisting}[caption={Observation construction.},label={lst:observation}]
def observation(self, agent_idx):
    obs = []
    # Self state
    obs.extend(velocity[agent_idx])      # 2 values
    obs.extend(position[agent_idx])      # 2 values

    # Agent ID one-hot
    one_hot = np.zeros(num_rescuers)
    one_hot[agent_idx] = 1.0
    obs.extend(one_hot)                  # num_rescuers values

    # Safe zones (relative pos + type)
    for zone in safe_zones:
        rel_pos = zone.pos - position[agent_idx]
        obs.extend(rel_pos)              # 2 values
        obs.append(zone.type / 3.0)      # 1 value (normalized)

    # Trees (relative pos, masked if not visible)
    for tree in trees:
        if is_visible(agent_idx, tree):
            obs.extend(tree.pos - position[agent_idx])
        else:
            obs.extend([0, 0])

    # Victims (relative pos + type, masked if not visible)
    for victim in victims:
        if is_visible(agent_idx, victim) and not saved[victim]:
            obs.extend(victim.pos - position[agent_idx])
            obs.append(victim.type / 3.0)
        else:
            obs.extend([0, 0, -1])

    return np.array(obs, dtype=np.float32)
\end{lstlisting}

% ----------------------------------------------------------------------------
\subsection{Training Pipeline: train.py}
\label{subsec:train_pipeline}
% ----------------------------------------------------------------------------

The \texttt{train.py} file implements the training process using Multi-Agent Proximal Policy Optimization (MAPPO) with TorchRL.
The pipeline is optimized for efficient reinforcement learning through parallel simulations, observation preprocessing, performance logging, and model-saving mechanisms.

\subsubsection{Reinforcement Learning Setup}
\label{subsubsec:rl_setup}

The MAPPO algorithm trains rescuer agents to maximize the cumulative reward function.
The algorithm optimizes the clipped surrogate objective as shown in Equation~\ref{eq:ppo_objective}.
\begin{equation}
\label{eq:ppo_objective}
L^{\text{PPO}}(\theta) = \mathbb{E}_t \left[ \min \left( r_t(\theta) \hat{A}_t, \text{clip}(r_t(\theta), 1-\epsilon, 1+\epsilon) \hat{A}_t \right) \right]
\end{equation}
where:
\begin{itemize}
    \item $r_t(\theta) = \frac{\pi_\theta(a_t|s_t)}{\pi_{\theta_{\text{old}}}(a_t|s_t)}$ is the probability ratio between current and old policies.
    \item $\hat{A}_t$ is the estimated advantage at time step $t$.
    \item $\epsilon = 0.2$ is the clipping threshold that prevents excessively large policy updates.
\end{itemize}

\subsubsection{Hyperparameter Configuration}
\label{subsubsec:hyperparams}

Key hyperparameters are carefully configured to optimize learning efficiency.

\begin{table}[H]
\centering
\begin{tabular}{lll}
\toprule
\textbf{Parameter} & \textbf{Default} & \textbf{Config} \\
\midrule
Learning rate $\alpha$ & $3 \times 10^{-4}$ & $3 \times 10^{-3}$ \\
Batch size & 256 & 2048 \\
Frames per batch & 2048 & 2048 \\
Number of epochs & 10 & 40 \\
Total timesteps & 100,000 & 204,800 \\
GAE $\gamma$ & 0.99 & 0.99 \\
GAE $\lambda$ & 0.95 & 0.95 \\
Clip epsilon $\epsilon$ & 0.2 & 0.2 \\
Entropy coefficient & 0.1 & 0.1 \\
Gradient clip norm & 1.0 & 1.0 \\
\bottomrule
\end{tabular}
\caption{Training hyperparameters (Default = code default; Config = conf/config.yaml).}
\label{tab:hyperparams}
\end{table}

\subsubsection{Training Process}
\label{subsubsec:training_process}

The training process follows the steps outlined in Algorithm~\ref{alg:mappo}.

\begin{algorithm}[H]
\caption{MAPPO Training Loop}
\label{alg:mappo}
\begin{algorithmic}[1]
\State Initialize policy network $\pi_\theta$ and value network $V_\phi$
\State Initialize data collector and replay buffer
\For{iteration $= 1$ to total\_iterations}
    \State Collect frames\_per\_batch steps using current policy
    \State Compute advantages using GAE: $\hat{A}_t = \sum_{l=0}^{\infty} (\gamma\lambda)^l \delta_{t+l}$
    \For{epoch $= 1$ to num\_epochs}
        \For{minibatch in replay buffer}
            \State Compute policy ratio $r_t(\theta)$
            \State Compute clipped objective $L^{\text{PPO}}$
            \State Compute value loss $L^{\text{value}}$
            \State Compute entropy bonus $L^{\text{entropy}}$
            \State Update $\theta, \phi$ using Adam optimizer
        \EndFor
    \EndFor
    \State Log metrics to TensorBoard
    \State Update policy in collector
\EndFor
\State Save final checkpoint
\end{algorithmic}
\end{algorithm}

\subsubsection{Performance Logging}
\label{subsubsec:logging}

Training progress is logged to facilitate real-time monitoring and offline analysis:

\begin{itemize}
    \item \textbf{TensorBoard Logging}: Metrics such as rewards, losses, and learning rates are visualized in real time.

    \item \textbf{Logged Metrics}:
    \begin{itemize}
        \item \texttt{train/loss/objective}: PPO clipped objective loss.
        \item \texttt{train/loss/critic}: Value function loss.
        \item \texttt{train/loss/entropy}: Entropy bonus (exploration).
        \item \texttt{train/loss/total}: Combined loss.
        \item \texttt{train/episode\_reward}: Mean episode reward.
        \item \texttt{train/total\_frames}: Cumulative frames processed.
    \end{itemize}
\end{itemize}

\subsubsection{Model Saving}
\label{subsubsec:model_saving}

Trained models are saved with configuration for reproducibility.

\begin{lstlisting}[caption={Checkpoint saving.},label={lst:checkpoint}]
checkpoint = {
    "policy_state_dict": policy.state_dict(),
    "critic_state_dict": critic.state_dict(),
    "optimizer_state_dict": optimizer.state_dict(),
    "total_frames": total_frames,
    "iteration": iteration,
    "env_config": env_kwargs
}
torch.save(checkpoint, f"{save_folder}/checkpoint.pt")
\end{lstlisting}

% ----------------------------------------------------------------------------
\subsection{Evaluation Pipeline: eval.py}
\label{subsec:eval_pipeline}
% ----------------------------------------------------------------------------

The \texttt{eval.py} script evaluates trained policies by running multiple episodes and computing performance metrics.

\subsubsection{Policy Loading}
\label{subsubsec:policy_loading}

The evaluation pipeline supports both manual and automatic policy loading.

\textbf{Manual Loading}: User provides the file path for the saved model, as shown in Equation~\ref{eq:manual_load}.
\begin{equation}
\label{eq:manual_load}
\pi_{\theta_{\text{loaded}}} = \text{load}(\text{path\_to\_model})
\end{equation}

\textbf{Automatic Loading}: The script identifies the latest saved model using file modification timestamps, as shown in Equation~\ref{eq:auto_load}.
\begin{equation}
\label{eq:auto_load}
\pi_{\theta_{\text{loaded}}} = \arg\max_\theta(t_{\text{saved}})
\end{equation}

\subsubsection{Evaluation Process}
\label{subsubsec:eval_process}

The evaluation runs the environment for $N_{\text{episodes}}$ episodes, where agents act based on the loaded policy $\pi_{\theta}$.
The environment produces a trajectory as shown in Equation~\ref{eq:trajectory}.
\begin{equation}
\label{eq:trajectory}
\tau = \{(s_t, a_t, r_t, s_{t+1}) \mid t = 0, 1, \ldots, T\}
\end{equation}

The cumulative reward for an episode is given by Equation~\ref{eq:episode_reward}.
\begin{equation}
\label{eq:episode_reward}
R_{\text{episode}} = \sum_{t=0}^{T} r_t
\end{equation}

\subsubsection{Performance Metrics}
\label{subsubsec:perf_metrics}

The following metrics are computed over evaluation episodes.

\textbf{Average Reward} is given by Equation~\ref{eq:avg_reward}.
\begin{equation}
\label{eq:avg_reward}
\bar{R} = \frac{1}{N_{\text{episodes}}} \sum_{i=1}^{N_{\text{episodes}}} R_{\text{episode},i}
\end{equation}

\textbf{Per-Agent Rewards}: Individual agent rewards $R_{\text{agent},i}$ are computed for each rescuer.

\textbf{Success Rate} is given by Equation~\ref{eq:success_rate}.
\begin{equation}
\label{eq:success_rate}
\text{Success Rate} = \frac{N_{\text{successful\_rescues}}}{N_{\text{victims}}}
\end{equation}

\subsubsection{Rendering}
\label{subsubsec:rendering}

To visually verify agent behavior, the evaluation pipeline supports rendering the environment in real time.

The rendering displays:
\begin{itemize}
    \item Safe zones with type-specific colors and transparency.
    \item Trees as gray circular obstacles.
    \item Victims with type-matching colors.
    \item Rescuers as white circles with vision radius indicators.
    \item Type labels (A, B, C, D) on zones and victims.
\end{itemize}

% ----------------------------------------------------------------------------
\subsection{Workflow Orchestration: main.py}
\label{subsec:main}
% ----------------------------------------------------------------------------

The \texttt{main.py} script serves as the central control hub for the project, orchestrating training and evaluation workflows using Hydra for configuration management.

\subsubsection{Core Functionality}
\label{subsubsec:core_func}

The script provides three primary execution modes:

\begin{enumerate}
    \item \textbf{Training Mode}: Invokes the \texttt{train()} function to train agents using MAPPO.
    \item \textbf{Evaluation Mode}: Invokes the \texttt{evaluate()} function to assess trained policies.
    \item \textbf{TensorBoard Mode}: Launches TensorBoard server for visualization.
\end{enumerate}

\subsubsection{Configuration Management}
\label{subsubsec:config_mgmt}

Hydra enables flexible configuration through YAML files and command-line overrides.

\begin{lstlisting}[language=bash,caption={Example usage.},label={lst:usage}]
# Training with default config
python -m src.main train.active=true

# Training with custom parameters
python -m src.main train.active=true \
    env.victims=12 env.rescuers=6 \
    train.total_timesteps=500000

# Evaluation with rendering
python -m src.main eval.active=true eval.render_mode=human
\end{lstlisting}

\subsubsection{Modular Design}
\label{subsubsec:modular}

The script is designed to be modular, enabling users to:
\begin{itemize}
    \item Easily switch between training and evaluation modes.
    \item Configure environment parameters through configuration files.
    \item Extend functionality for fine-tuning, testing, or benchmarking policies.
\end{itemize}

% ============================================================================
\section{Environment Parameters}
\label{sec:parameters}
% ============================================================================

This section provides a comprehensive reference of all configurable environment parameters.

\begin{table}[H]
\centering
\begin{tabular}{llll}
\toprule
\textbf{Parameter} & \textbf{Default} & \textbf{Config} & \textbf{Description} \\
\midrule
\texttt{num\_rescuers} & 2 & 6 & Number of rescuer agents \\
\texttt{num\_victims} & 2 & 12 & Number of victim entities \\
\texttt{num\_trees} & 5 & 8 & Number of obstacle trees \\
\texttt{num\_safe\_zones} & 4 & 4 & Number of safe zones (corners) \\
\texttt{max\_cycles} & 200 & 300 & Maximum steps per episode \\
\texttt{vision\_radius} & 0.5 & 0.4 & Maximum observation distance \\
\texttt{rescue\_radius} & 0.15 & 0.15 & Distance for successful rescue \\
\texttt{follow\_radius} & 0.2 & 0.2 & Distance for victim commitment \\
\texttt{agent\_size} & 0.03 & 0.03 & Radius of agent entities \\
\texttt{tree\_radius} & 0.05 & 0.05 & Radius of tree obstacles \\
\texttt{safe\_zone\_radius} & 0.15 & 0.15 & Radius of safe zones \\
\texttt{world\_size} & 2.0 & 2.0 & World bounds ([-1, 1] range) \\
\texttt{continuous\_actions} & True & True & Use continuous action space \\
\texttt{randomize\_safe\_zones} & False & True & Randomize safe zone types \\
\bottomrule
\end{tabular}
\caption{Environment parameters (Default = code default; Config = conf/config.yaml).}
\label{tab:env_params}
\end{table}

% ============================================================================
\section{Model Architecture}
\label{sec:model_architecture}
% ============================================================================

The neural network architecture follows the MAPPO paradigm with decentralized actors and centralized critics.

% ----------------------------------------------------------------------------
\subsection{Policy Network (Actor)}
\label{subsec:policy_network}
% ----------------------------------------------------------------------------

The policy network maps local observations to action distributions:

\begin{itemize}
    \item \textbf{Input}: Local observation $[\text{batch}, N_{\text{agents}}, D_{\text{obs}}]$.
    \item \textbf{Architecture}: 2-layer MLP with 64 hidden units per layer.
    \item \textbf{Activation}: Tanh activation function.
    \item \textbf{Output}: Parameters for TanhNormal distribution (location and scale).
    \item \textbf{Parameter Sharing}: Shared across all agents (homogeneous policy).
\end{itemize}

The policy architecture can be expressed as shown in Equations~\ref{eq:policy_h1}--\ref{eq:policy_sigma}.
\begin{align}
    \mathbf{h}_1 &= \tanh(W_1 \mathbf{o} + b_1) \label{eq:policy_h1} \\
    \mathbf{h}_2 &= \tanh(W_2 \mathbf{h}_1 + b_2) \label{eq:policy_h2} \\
    \boldsymbol{\mu} &= W_\mu \mathbf{h}_2 + b_\mu \label{eq:policy_mu} \\
    \boldsymbol{\sigma} &= \text{softplus}(W_\sigma \mathbf{h}_2 + b_\sigma) \label{eq:policy_sigma}
\end{align}

Actions are sampled from: $\mathbf{a} \sim \text{TanhNormal}(\boldsymbol{\mu}, \boldsymbol{\sigma})$.

% ----------------------------------------------------------------------------
\subsection{Value Network (Critic)}
\label{subsec:value_network}
% ----------------------------------------------------------------------------

The value network estimates state values for advantage computation:

\begin{itemize}
    \item \textbf{Input}: All agent observations $[\text{batch}, N_{\text{agents}}, D_{\text{obs}}]$.
    \item \textbf{Architecture}: 2-layer MLP with 128 hidden units per layer.
    \item \textbf{Activation}: Tanh activation function.
    \item \textbf{Output}: State value $[\text{batch}, N_{\text{agents}}, 1]$.
    \item \textbf{Centralized}: True (observes all agent observations for credit assignment).
\end{itemize}

% ============================================================================
\section{Curriculum Learning}
\label{sec:curriculum}
% ============================================================================

The project supports curriculum learning for progressive difficulty adjustment during training.

\subsection{Curriculum Configuration}
\label{subsec:curriculum_config}

\begin{itemize}
    \item \textbf{Minimum Trees}: Starting number of obstacles (default: 0).
    \item \textbf{Maximum Trees}: Final number of obstacles (default: 8).
    \item \textbf{Number of Stages}: Progression stages (default: 5).
    \item \textbf{Stage Duration}: Automatically calculated as $\frac{\text{total\_iterations}}{\text{num\_stages}}$.
\end{itemize}

\subsection{Progression Strategy}
\label{subsec:progression}

The number of trees at stage $s$ is computed as shown in Equation~\ref{eq:curriculum}.
\begin{equation}
\label{eq:curriculum}
N_{\text{trees}}(s) = N_{\min} + \left\lfloor \frac{s}{S-1} \times (N_{\max} - N_{\min}) \right\rfloor
\end{equation}
where $S$ is the total number of stages.

The curriculum maintains a fixed observation space size by using \texttt{max\_trees} for observation dimensionality, with unused tree slots masked as zeros.

% ============================================================================
\section{Usage Guide}
\label{sec:usage}
% ============================================================================

\subsection{Docker Execution}
\label{subsec:docker}

\begin{lstlisting}[language=bash,caption={Docker execution commands.},label={lst:docker}]
# Build image
docker build -t student-search:latest -f docker/Dockerfile .

# Training
docker run --rm \
    -v "${PWD}/search_rescue_logs:/app/search_rescue_logs" \
    student-search:latest

# Evaluation
docker run --rm \
    -v "${PWD}/search_rescue_logs:/app/search_rescue_logs" \
    student-search:latest eval.active=true
\end{lstlisting}

\subsection{Local Execution}
\label{subsec:local}

\begin{lstlisting}[language=bash,caption={Local execution commands.},label={lst:local}]
# Install dependencies
pip install -e .
pip install -e ".[dev]"

# Training
python -m src.main train.active=true

# Evaluation with rendering
python -m src.main eval.active=true eval.render_mode=human

# Custom configuration
python -m src.main train.active=true \
    env.victims=12 env.rescuers=6 \
    train.total_timesteps=500000 \
    train.learning_rate=0.001
\end{lstlisting}

% ============================================================================
\section{Conclusion}
\label{sec:conclusion}
% ============================================================================

This documentation has provided a comprehensive overview of the Search and Rescue Multi-Agent Simulation project.
The key contributions include:

\begin{enumerate}
    \item A flexible multi-agent environment implementing realistic search and rescue dynamics.
    \item A commitment-based victim following system with hysteresis for stable agent-victim relationships.
    \item A carefully designed reward structure that enables effective credit assignment.
    \item Integration with TorchRL for efficient MAPPO training.
    \item Optional curriculum learning for progressive difficulty adjustment.
    \item Comprehensive logging and visualization capabilities.
\end{enumerate}

The modular design allows for easy extension and experimentation with different algorithms, environment configurations, and training strategies.

\end{document}
